% =============================================================================
\chapter{Cronograma Preliminar}
% =============================================================================

\section{Fase 1 --- fundamentação e posicionamento}

\begin{itemize}
\item concluir revisão estruturada de HTN, MORL e arquiteturas híbridas;
\item consolidar lacuna, questão de pesquisa e hipóteses;
\item revisar a especificação do domínio experimental.
\end{itemize}

\section{Fase 2 --- formalização}

\begin{itemize}
\item formalizar estado, tarefas, métodos, preferências e recompensas;
\item definir a unidade temporal de cada decisão MORL;
\item especificar as interfaces e os registros de explicação.
\end{itemize}

\section{Fase 3 --- estabilização do baseline}

\begin{itemize}
\item concluir testes do HTN, sensores, execução e replanejamento;
\item desacoplar seleção de método da ordem fixa;
\item implementar máscara e identificadores estáveis de métodos.
\end{itemize}

\section{Fase 4 --- ambiente multiobjetivo}

\begin{itemize}
\item implementar o cenário tático ou de sobrevivência;
\item definir eventos, recompensas vetoriais e sementes;
\item implementar instrumentação e execução automatizada.
\end{itemize}

\section{Fase 5 --- preferências}

\begin{itemize}
\item implementar o GOWG determinístico;
\item validar normalização e mudanças contextuais;
\item criar perfis e condições de teste.
\end{itemize}

\section{Fase 6 --- MORL e integração}

\begin{itemize}
\item selecionar e implementar o algoritmo MORL;
\item integrar máscara de métodos válidos;
\item treinar a política sob múltiplas preferências;
\item validar o ciclo online HTN--GOWG--MORL.
\end{itemize}

\section{Fase 7 --- baselines e ablações}

\begin{itemize}
\item implementar os quatro baselines;
\item implementar configurações de ablação;
\item validar equivalência de cenários e orçamento.
\end{itemize}

\section{Fase 8 --- avaliação}

\begin{itemize}
\item executar experimentos com múltiplas sementes;
\item testar mudanças de preferência e generalização;
\item analisar desempenho, validade, diversidade e custo;
\item produzir tabelas, curvas e traços de decisão.
\end{itemize}

\section{Fase 9 --- dissertação e publicação}

\begin{itemize}
\item consolidar metodologia e resultados;
\item discutir limitações e ameaças à validade;
\item redigir a dissertação e um artigo científico;
\item preparar artefatos de reprodutibilidade.
\end{itemize}
